% Generated draft: requires mathematical and source review.
\documentclass[11pt]{article}
\usepackage[margin=25mm]{geometry}
\usepackage{fontspec}
\setmainfont{Latin Modern Roman}
\usepackage{amsmath,amssymb,mathtools,unicode-math}
\setmathfont{Latin Modern Math}
\usepackage{enumitem,needspace,xurl,hyperref}
\setlength{\parindent}{0pt}
\setlength{\parskip}{5pt}
\setlength{\emergencystretch}{3em}
\newcommand{\N}{\mathbb{N}}
\newcommand{\Z}{\mathbb{Z}}
\newcommand{\Q}{\mathbb{Q}}
\newcommand{\R}{\mathbb{R}}
\newcommand{\C}{\mathbb{C}}
\newcommand{\F}{\mathbb{F}}
\newcommand{\A}{\mathbb{A}}
\newcommand{\PP}{\mathbb{P}}
\newcommand{\E}{\mathbb{E}}
\newcommand{\eps}{\varepsilon}
\newcommand{\dd}{\,\mathrm{d}}
\begin{document}
\section*{UnsolvedMath ID: 30001581}
\textbf{Source problem code:} catalogue entry
\subsection{Definitions and mathematical statement}
Let \(K\) be an imaginary quadratic field, let \(B\) be a quaternion division algebra over \(K\), and let \(M=\Gamma\backslash\mathbb H^3\) be an arithmetic hyperbolic three-manifold associated with \(B\), as specified in the supplied source. Here \(\Gamma\) is a torsion-free arithmetic lattice and \(\mathbb H^3\) has sectional curvature \(-1\); volume is measured in this normalization.

Let \((M_N\to M)_{N\geq 1}\) be a residual chain of finite-sheeted congruence covers. In subgroup notation, write
\[
M_N=\Gamma_N\backslash\mathbb H^3,\qquad
\Gamma\supseteq\Gamma_1\supseteq\Gamma_2\supseteq\cdots,\qquad
[\Gamma:\Gamma_N]<\infty,\qquad
\bigcap_{N\geq1}\Gamma_N=\{1\},
\]
where each \(\Gamma_N\) is a congruence subgroup. Congruence refers to containing a kernel of reduction modulo a nonzero ideal in an integral arithmetic model for the lattice. The supplied record does not specify that model or whether the term ``residual chain'' additionally requires the subgroups to be normal in \(\Gamma\). Any such intended requirement remains part of the unresolved source convention; the displayed conditions do not resolve it.

For an abelian group \(A\), define
\[
A_{\mathrm{tors}}=\{a\in A: ma=0\text{ for some integer }m\geq1\}.
\]
Thus \(H_1(M_N,\mathbb Z)_{\mathrm{tors}}\) is the finite torsion subgroup of the first integral homology group. Write \(\#\) for cardinality and \(\log\) for the natural logarithm.

\textbf{Conjectural statement.} For every manifold \(M\) of the specified arithmetic type and every residual chain of congruence covers in the source's intended sense, is it true that
\[
\lim_{N\to\infty}
\frac{\log\!\left(\#H_1(M_N,\mathbb Z)_{\mathrm{tors}}\right)}
{\operatorname{vol}(M_N)}
=\frac{1}{6\pi}\,?
\]
The quantifier concerns congruence covers of the specified arithmetic manifolds. The broader discussion of arbitrary residual chains in the supplied literature assessment does not enlarge this formulation. The source convention concerning normality should be checked before treating the quantified class of chains as fully resolved.
\subsection{Short English statement}
For an arithmetic hyperbolic three-manifold arising from a quaternion division algebra over an imaginary quadratic field, the conjecture predicts that along every residual chain of congruence covers, the logarithm of the size of the torsion in first integral homology, divided by hyperbolic volume, tends to \(1/(6\pi)\). The supplied record labels the problem open; the precise convention for residual chains requires review.
\subsection{Sources}
\begin{itemize}
\item Supplied local UnsolvedMath registry record, permanent ID \(30001581\), source problem code \texttt{catalogue entry}, title ``Homological Torsion Growth in Arithmetic Hyperbolic Towers,'' in the Oberwolfach Reports Open Problems collection. No individual proposer is supplied; the proposed year is recorded as 2010.
\item Supplied source citation: \emph{Actions and Invariants of Residually Finite Groups: Asymptotic Methods} (2011), Oberwolfach Reports, pp.~2335--2391, DOI \texttt{10.4171/owr/2010/40}. Original OWR report: \url{https://doi.org/10.4171/owr/2010/40}. The supplied record attributes the problem to workshop open-problem discussions documented by this report.
\item Supplied literature reference: Bergeron--Venkatesh, ``torsion homology growth,'' \url{https://arxiv.org/abs/1004.1083}.
\item Literature assessment checked at \texttt{2026-08-21}: this is supplied dated curation, not independently verified here. It labels the problem open and states that the residual-chain torsion-growth conjecture remains open in general, while results for cyclic, congruence, and selected arithmetic towers do not settle arbitrary residual chains of hyperbolic three-manifolds. The supplied background explicitly requests independent verification of time-sensitive status and cited claims.
\item The registry labels the statement \texttt{corrected\_verified} and reports that its mathematical content was checked against the official Oberwolfach report, with notation, context, and question form normalized. This is a supplied verification claim, not a source check performed for this draft.
\end{itemize}
\end{document}
